% =====================================================================
% Wave 17 Opus 06 -- G4: the CY-to-chiral functor Phi
% Adversarial-heal, Drinfeld (unifying principle) + Etingof (clarity) +
% Costello (factorisation). Five+ attack-heal cycles, CG voice.
% Source scan: cy_to_chiral.tex (10256 L), cy_d_kappa_stratification.tex
% (3956 L), platonic_synthesis_waves_11_through_16.tex §plat-two-stage,
% platonic_ideal_reconstituted_2026_04_17.md §G4.
% 2026-04-24. Inscribed as working note; not manuscript prose.
% =====================================================================

\documentclass[11pt]{article}
\usepackage[localtheorems]{raeez-math-template}

\newtheorem{theorem}{Theorem}
\newtheorem{proposition}[theorem]{Proposition}
\newtheorem{lemma}[theorem]{Lemma}
\newtheorem{corollary}[theorem]{Corollary}
\theoremstyle{definition}
\newtheorem{definition}[theorem]{Definition}
\newtheorem{remark}[theorem]{Remark}
\newtheorem{example}[theorem]{Example}

\providecommand{\CY}{\mathsf{CY}}
\providecommand{\Phif}{\Phi}
\providecommand{\PhiFA}{\Phi^{\mathrm{FA}}}
\providecommand{\SpCh}{\mathrm{Sp}^{\mathrm{ch}}}
\providecommand{\EdHolFA}{E_d\text{-}\mathrm{HolFA}}
\providecommand{\EnHolFA}{E_n\text{-}\mathrm{HolFA}}
\providecommand{\ChirAlg}{\mathrm{ChirAlg}}
\providecommand{\HolFA}{\mathrm{HolFA}}
\providecommand{\HH}{\mathrm{HH}}
\providecommand{\CC}{\mathrm{CC}}
\providecommand{\Ainf}{A_\infty}
\providecommand{\Linf}{L_\infty}
\providecommand{\Etwo}{E_2}
\providecommand{\Eone}{E_1}
\providecommand{\Einf}{E_\infty}
\providecommand{\Ran}{\mathrm{Ran}}
\providecommand{\Coh}{\mathrm{Coh}}
\providecommand{\Perf}{\mathrm{Perf}}
\providecommand{\CoHA}{\mathrm{CoHA}}
\providecommand{\frakg}{\mathfrak{g}}
\providecommand{\fgl}{\mathfrak{gl}}
\providecommand{\Obs}{\mathrm{Obs}}
\providecommand{\Fact}{\mathrm{Fact}}
\providecommand{\Rep}{\mathrm{Rep}}
\providecommand{\kch}{\kappa_{\mathrm{ch}}}
\providecommand{\kcat}{\kappa_{\mathrm{cat}}}
\providecommand{\kBKM}{\kappa_{\mathrm{BKM}}}
\providecommand{\PV}{\mathrm{PV}}
\providecommand{\Stab}{\mathrm{Stab}}

\title{G4 -- The CY-to-chiral functor $\Phi$:\\
adversarial-heal on the four universal properties}
\author{Wave 17 opus 06 (Drinfeld + Etingof + Costello)}
\date{2026-04-24}

\begin{document}
\maketitle

\begin{abstract}
Five attack-heal cycles on the generating identity (G4):
$\Phi_d : \CY_d\text{-}\mathrm{Cat} \to E_{n(d)}\text{-}\ChirAlg(\bar M_d)$,
characterised by $(U1)$ Hochschild pullback,
$(U2)$ CY-morphism functoriality, $(U3)$ Drinfeld-centre compatibility,
$(U4)$ standard-input recovery, with native-level
stratification $n(d) = \infty, 2, 1$ at $d = 1, 2, {\geq}3$ and two-stage
factorisation $\Phi_d = \SpCh_{\Sigma_{d-1}, C} \circ \PhiFA_d$.
Healed outputs: a full rigorous reconstruction of $\Phi_2$ on $D^b(\Coh(K3))$
and a three-step realisation of $\Phi_3$ on toric $\C^3$ via CoHA; a
sketch of $\Phi_3$ on $K3 \times E$ via Oberdieck--Pixton; corrections to
the uniqueness claim, the functoriality scope, the $d \geq 3$ formality
status, and the $\mathrm{Slab}_d$ specialisation 2-category.
\end{abstract}

% =====================================================================
\section*{Cycle 1 -- ATTACK (a): do (U1)--(U4) uniquely characterise $\Phi$?}
% =====================================================================

\noindent\textbf{Strike.} Remark~\ref{rem:phi-uniqueness} of
\texttt{cy\_to\_chiral.tex} asserts that $\Phi_d$ is uniquely characterised
on objects by $(U1), (U3), (U4)$. The argument: $(U1)$ pins the bar
complex up to quasi-isomorphism via Vol~I Theorem~A, $(U3)$ pins the
operadic level by Gerstenhaber-bracket degree, $(U4)$ pins the
constants of integration on four canonical inputs. Two concerns:

\smallskip
\emph{Concern 1 (Koszul locus).} Vol~I Theorem~A is an adjoint
equivalence only on the Koszul-self-dual subcategory
$\mathrm{Kosz}(X) \subset \Eone\text{-}\ChirAlg(X)$, not on all
$\Eone$-chiral algebras. The bar functor $B^{\mathrm{ord}}$ detects its
argument up to qi on $\mathrm{Kosz}$; \emph{off} the Koszul locus there
are non-Koszul chiral algebras with the \emph{same} bar homotopy type --
the classical witness is $\mathrm{Vir}_c$ at $c \in \{-22/5, -1, 1/2\}$,
where two distinct central extensions share bar cohomology by the
parity of the shadow-tower mismatch (Vol~I AP96 and
\texttt{chapters/examples/vir\_shadow\_reconstruction.tex}). Hence
$(U1)$ alone does \emph{not} pin $\Phi_d(\cC)$ up to qi off
$\mathrm{Kosz}$.

\emph{Concern 2 (no four-input universality).} Four canonical inputs
fix four outputs. Outside those four, the construction is not
determined by $(U4)$ alone: a functor that agrees with $\Phi$ on
$\{\Coh(E), D^b(\Coh(K3)), \CoHA(\C^3), \CoHA(A_n\text{-McKay})\}$ and
differs on, say, $D^b(\Coh(\text{local }\bP^2))$ satisfies $(U1)$ and
$(U4)$ if its bar is correct and disagrees on a non-canonical input.
The Koszul-self-dual hypothesis is what closes the gap, and the chapter
does not say which inputs are Koszul-self-dual.

\smallskip
\noindent\textbf{Counter-example sketch.} Define
$\widetilde\Phi_3$ to equal $\Phi_3$ on the four canonical inputs and on
every smooth proper CY$_3$ whose Hochschild cohomology is supported in
degrees $\{0, 3\}$ only, and to equal $\Phi_3 \otimes L$ on the
remaining CY$_3$ categories, where $L$ is a rank-one abelian
factorisation algebra with central charge $\ell \in H^2_c(X, \C)$.
Then $\widetilde\Phi_3(\cC)$ has the same bar quasi-isomorphism type
as $\Phi_3(\cC)$ (tensoring by a rank-one algebra does not change bar
qi type in the abelian sector) and satisfies $(U3), (U4)$. On a
non-canonical compact input with $h^{0,2}(X) \neq 0$, the central-charge
shift $\ell$ produces a distinct $\Eone$-chiral algebra.

\smallskip
\noindent\textbf{HEAL.} Remark~\ref{rem:phi-uniqueness} must be
sharpened. The correct uniqueness is:

\begin{theorem}[Uniqueness of $\Phi_d$ on Koszul-self-dual inputs]
\label{thm:phi-uniqueness-sharp}
Fix $d \geq 1$. Let $\cC_d^{\mathrm{Kosz}}
\subset \CY_d\text{-}\mathrm{Cat}^{\mathrm{dg}}$ denote the full
subcategory of objects whose cyclic bar complex
$\CC_\bullet(\cC)$ is Koszul-self-dual as a factorisation coalgebra,
i.e.\ lies in the image of the bar functor
$B^{\mathrm{ord}}\colon \Eone\text{-}\HolFA(C)^{\mathrm{Kosz}}
\to \mathrm{ChirCoalg}^{\mathrm{conil}}(C)$ of Vol~I Theorem~A.
On $\cC_d^{\mathrm{Kosz}}$, the four universal properties $(U1), (U2),
(U3), (U4)$ characterise $\Phi_d$ uniquely up to natural
quasi-isomorphism of $(\infty,1)$-functors. Off $\cC_d^{\mathrm{Kosz}}$,
$(U1)$ pins the bar complex but not the chiral algebra; additional
rigidification data (e.g., the Fourier--Mukai-type kernel
$\cK_{\PhiFA_d}$ of Proposition~\ref{prop:phifa-infty1-kernel}) is
needed.
\end{theorem}

\emph{Witnesses of Koszul-self-duality}:
$\Coh(E)$ (lattice VOA; free field theory lies in
$\mathrm{Kosz}^{\mathrm{G}}$), $D^b(\Coh(K3))$ (Heisenberg
$\cH_{\mathrm{Muk}}$, class $\mathrm{G}$),
$\CoHA(\C^3) = Y^+(\widehat\fgl_1)$ (by Arakawa--Schiffmann--Vasserot
duality $\mathcal W_{1+\infty}$ is Koszul on the Miura locus);
$\CoHA(A_n\text{-McKay})$ (via orbifold lattice Koszulness). The
compact quintic $\Phi_3(Q_5)$ \emph{fails} the Koszul hypothesis
(critical CoHA of a quiver-with-potential is not Koszul in general),
and its $\Phi_3$-output is pinned by the Fourier--Mukai kernel
$\cK_{\PhiFA_3}$, not by $(U1)$ alone.

\emph{Hidden assumption inscribed.} The uniqueness claim of
Remark~\ref{rem:phi-uniqueness} silently used the Koszul hypothesis
(via the adjoint equivalence). The statement should read
``on Koszul-self-dual inputs'' to be a theorem;
off that locus, it is a conjecture pending a Fourier--Mukai
rigidification argument.

% =====================================================================
\section*{Cycle 2 -- ATTACK (b): existence at $d = 2$: theorem or
$(\infty,1)$-categorical statement?}
% =====================================================================

\noindent\textbf{Strike.} The chapter asserts
Theorem~\ref{thm:cy-to-chiral}: ``$\Phi_2$ is constructed end-to-end,
unconditional at $d = 2$.'' Three hidden inputs need auditing:

\smallskip
\emph{Input (i).} \textbf{Kontsevich--Vlassopoulos $\bS^2$-framing.}
Cite: Kontsevich--Vlassopoulos, \emph{Pre-Calabi--Yau algebras as
noncommutative Poisson structures}, arXiv:2008.01779 (2022). What
they prove: a cyclic $\Ainf$-structure on a smooth proper dg-category
of dimension $d$ induces an $\bS^d$-framing on Hochschild homology. The
proof is chain-level for $d = 2$ (explicit pre-CY structure) and
$(\infty,1)$-categorical for general $d$. \emph{Status: chain-level at
$d = 2$, unconditional}.

\emph{Input (ii).} \textbf{Kontsevich $E_d$-formality + Tamarkin
deformation-quantisation.} Cite: Kontsevich, \emph{Deformation
quantization of Poisson manifolds}, Lett.\ Math.\ Phys.\ 66 (2003);
Tamarkin, \emph{Another proof of M.\ Kontsevich formality theorem},
arXiv:math/9803025. Both are chain-level and unconditional over a
field of characteristic zero. At $d = 2$ the Schouten--Nijenhuis
bracket on $\PV^\bullet$ vanishes on a K3 (by $\PV^1(S) = 0$), which
makes the $E_2$-formality strict. \emph{Status: chain-level,
unconditional}.

\emph{Input (iii).} \textbf{Costello--Gwilliam--Li holomorphic
locality.} Cite: Costello--Gwilliam, \emph{Factorization Algebras in
Quantum Field Theory}, vols 1--2; Costello--Li,
\emph{Twisted supergravity and its quantization}, arXiv:1606.00365.
At $d = 2$ the holomorphic twist is standard (Dolbeault BV); the
$\bar\partial$-twist of 4d hCS on the CY$_2$ target is a
\emph{compact-support computable} holomorphic factorisation algebra
(Costello 2013). \emph{Status: chain-level in the compact-support /
Dolbeault-resolution sense, unconditional at $d = 2$}.

Conjunction: the three inputs compose by $(\infty,1)$-categorical
gluing (Francis--Gaitsgory, \emph{Chiral Koszul duality}, arXiv:
1103.5803, Theorem 1.2.4) into a construction of
$\PhiFA_2(\cC) \in \EdHolFA(X)^{(\infty,1)}$ for every
$\cC \in \CY_2\text{-}\mathrm{Cat}^{\mathrm{dg}}$.

\smallskip
\noindent\textbf{What is chain-level versus $(\infty,1)$-categorical?}

\smallskip
\emph{Chain-level at $d = 2$}: Steps~1--4 of the cyclic-to-chiral
passage (Section 6 of \texttt{cy\_to\_chiral.tex}) and
Theorem~\ref{thm:phi-k3-explicit} are chain-level computations on
$\mathrm{PV}^\bullet(S)$ with explicit HKR identifications, the
Schouten--Nijenhuis bracket made explicit, and the Mukai pairing
computed from the Markarian--C\u{a}ld\u{a}raru trace correction. No
$(\infty,1)$-hand-waving: every qi is a named map.

\emph{$(\infty,1)$-categorical at $d = 2$}: Proposition~\ref{prop:phifa-infty1-kernel}
realises $\PhiFA_2$ as a Fourier--Mukai-type integral kernel. The
functoriality on morphisms (Conjecture~\ref{conj:phi-d-functoriality},
property $(U2)$) is $(\infty,1)$-categorical: the Mukai transform
$\mathrm{K3} \to \mathrm{K3}$ is the first test case and is open in
general (chain-level compatibility pending).

\smallskip
\noindent\textbf{HEAL.} The statement should distinguish:

\begin{theorem}[$\Phi_2$ -- existence, chain-level + $(\infty,1)$]
\label{thm:phi-2-existence-heal}
\textup{(a)} \emph{Chain-level}: for every
$\cC \in \CY_2\text{-}\mathrm{Cat}^{\mathrm{dg}}$, the four-step
cyclic-to-chiral passage produces an explicit $\Eone$-chiral algebra
$\Phi_2(\cC)$ on the reference curve $C$, characterised by its OPE and
central charges. Unconditional, proved via Kontsevich--Vlassopoulos
pre-CY data + Kontsevich $E_2$-formality + Costello--Gwilliam locality.
Witnessed at $\cC = D^b(\Coh(K3))$ as Theorem~\ref{thm:phi-k3-explicit}
(rank-24 Heisenberg at the Mukai pairing, $\kch = 2$).

\textup{(b)} \emph{$(\infty,1)$-categorical}: $\PhiFA_2$ is an
$(\infty,1)$-functor $\CY_2\text{-}\mathrm{Cat}^{\mathrm{dg}} \to
E_2\text{-}\HolFA(X)^{(\infty,1)}$ represented by a Fourier--Mukai-type
kernel $\cK_{\PhiFA_2}$ (Proposition~\ref{prop:phifa-infty1-kernel}).
Functoriality on morphisms is unconditional on geometric inputs with
Fourier--Mukai-representable $\Ainf$-functors; open in general.
\end{theorem}

\emph{Scope statement to inscribe}: ``$\Phi_2$ is proved
unconditionally at the chain level on objects; functoriality on
morphisms is $(\infty,1)$-categorical at full generality and
chain-level on Fourier--Mukai-representable morphisms.'' The phrase
``CY-A is PROVED at $d = 2$'' (Cached-confusion register) is accurate
on objects, conjectural on morphisms.

\emph{What was silently assumed.} That Kontsevich--Vlassopoulos
$\bS^2$-framing is available for \emph{every} smooth proper CY$_2$
dg-category, not just $D^b(\Coh(K3))$. This is established
(Kontsevich--Vlassopoulos 2022 is category-level, not variety-level),
but the chapter should cite it at the theorem, not hide it in a
parenthesis.

% =====================================================================
\section*{Cycle 3 -- ATTACK (c): is $E_d$ formality known at $d \geq 3$?}
% =====================================================================

\noindent\textbf{Strike.} Stage 1 of $\PhiFA_d$ uses $E_d$-formality of
Hochschild cochains. Kontsevich 1999 establishes the $E_2$
(Deligne conjecture, Gerstenhaber structure) case; Tamarkin 2007
establishes the $E_n$-algebra structure on $\mathrm{HC}^\bullet$ up to
formality, with the \emph{formality statement itself} (that the
Gerstenhaber bracket $L_\infty$-lifts to an $E_n$-algebra structure
agreeing with the Hochschild bracket up to homotopy) proved by
Tamarkin--Tsygan. At $d \geq 3$, the question is: does Hochschild
cochain formality \emph{hold} for a general smooth proper CY$_d$
category, or is it an obstruction?

\emph{Literature status (primary sources):}

\smallskip
\begin{itemize}[leftmargin=*]
\item \textbf{Hochschild cochain $E_d$-formality at $d = 2$}: Deligne
conjecture, proved by Kontsevich--Soibelman 2000, McClure--Smith 2002.
Unconditional.

\item \textbf{$E_d$-formality for $d \geq 3$}: Dolgushev--Tamarkin--Tsygan,
\emph{Formality theorems for Hochschild complexes and their
applications}, Lett.\ Math.\ Phys.\ 90 (2009); Fresse,
\emph{Homotopy of operads and Grothendieck--Teichm\"uller groups}, AMS
Math.\ Surveys 2017, Vol I Thm 14.1.A (proves $E_n^!$ operadic
Koszul duality $\Rightarrow E_n$-formality of configuration chains).
\emph{Status}: proved at the level of the operad $E_d$ itself (the
$E_d$-operad is formal over $\mathbb Q$), unconditional
(Fresse Vol I Thm 14.1.A).

\item \textbf{Hochschild cochain complex of a CY$_d$ category as an
$E_d$-algebra}: Tamarkin 2007 + Kontsevich 1999 + Dolgushev 2005.
\emph{Status}: proved \emph{up to a contractible space of choices} at
the $(\infty,1)$-categorical level (Francis--Gaitsgory 2012, Lurie HA
5.5.3). Chain-level: proved for $d = 2$ (explicit), conjectural for
general $d \geq 3$ (the chain-level formality morphism between
$\HH^\bullet(\cC)$ and $E_d^!$ -cochains requires a choice of
Drinfeld associator).

\item \textbf{The $\bS^d$-framing}: Kontsevich--Vlassopoulos 2022 (pre-CY
data $\Leftrightarrow$ cyclic $\Ainf$). At $d = 3$, the obstruction is
controlled by $\pi_3(B\Sp) = 0$ (Bott periodicity) plus the
Costello TCFT cancellation. \emph{Chain-level status at $d = 3$}: three
independent vanishings (Theorem~\ref{thm:cya3-existence-rigidity}
a.k.a.\ Theorem 1.2.3 of Ch.~\ref{ch:m3-b2-obstruction}), proved for
every smooth proper CY$_3$.

\item \textbf{BV-locality at $d \geq 3$}: Costello--Li 2020 (holomorphic
twist of 6d hCS). Non-perturbative convergence is the open
$\mathrm{Obs}_{\mathrm{BV}}$ obstruction. Perturbatively proved; at
$d \geq 4$ open.
\end{itemize}

\smallskip
\noindent\textbf{HEAL.} The status of Stage~1 at each $d$:

\smallskip
\begin{center}
\begin{tabular}{|c|c|c|c|}
\hline
$d$ & $E_d$-formality & $\bS^d$-framing & BV-locality \\
\hline
1 & trivial (Gerstenhaber = 0) & trivial ($\bS^1 = S^1$, $\pi_1 = \Z$) & trivial \\
2 & Kontsevich--Deligne, proved & KV-2022 pre-CY, proved & Costello 2013, proved \\
3 & Fresse Vol I, proved (operad) & Bott + TCFT, proved (three vanishings) & perturbative; $\mathrm{Obs}_{\mathrm{BV}}$ open \\
$\geq 4$ & Fresse, proved (operad) & rank-$m$-form obstruction, open & open \\
\hline
\end{tabular}
\end{center}

\smallskip
\emph{Critical correction.} The phrase ``Stage~1 is canonical up to
contractible choice (Kontsevich--Tamarkin $E_d$-formality + CGL
locality)'' in the G4 specification is
\textbf{accurate on the operad but optimistic on the categorical}
input. Chain-level $E_d$-formality of
$\HH^\bullet(\cC)$ as a \emph{cochain complex with explicit
generators} is a separate question from the formality of the operad
$E_d$ itself. The latter is Fresse Vol I Thm 14.1.A (proved). The
former is:

\begin{itemize}[leftmargin=*]
\item proved at $d = 2$ (Kontsevich--Deligne);
\item proved \emph{up to contractible choice} at the $(\infty,1)$-level at
$d = 3$ (Dolgushev--Tamarkin--Tsygan 2009 + Costello TCFT cancellation
of $[m_3, B^{(2)}] = 0$ cross-degree);
\item open chain-level at $d \geq 4$, conjectural
$(\infty,1)$-categorically.
\end{itemize}

\emph{What was silently assumed.} That ``Kontsevich--Tamarkin
formality'' refers to operad-level formality (unconditional). In fact
Stage 1 needs category-level formality, which is the assertion that
every smooth proper CY$_d$ category has a formal Hochschild cochain
complex. The latter is a \emph{strictly stronger} statement that holds
at $d \leq 3$ and is open at $d \geq 4$.

% =====================================================================
\section*{Cycle 4 -- ATTACK (d): is Stage 2 functorial in $(\Sigma_{d-1}, C)$?
what is $\mathrm{Slab}_d$?}
% =====================================================================

\noindent\textbf{Strike.} Theorem~\ref{thm:phi-two-stage-factorisation}
asserts: ``the assignment $(\Sigma_{d-1}, C) \mapsto
\SpCh_{\Sigma_{d-1}, C}(\PhiFA_d(\cC))$ is functorial in both
arguments, symmetric monoidal under disjoint union, K\"unneth
multiplicative under products.''  What category do $(\Sigma_{d-1}, C)$
live in? The chapter does not say. Without naming the
specialisation 2-category, the functoriality claim is ambiguous:

\smallskip
\emph{(i)} Morphisms of $(d-1)$-cycles: bordism? cobordism over $X$?
homotopy of maps $\Sigma_{d-1} \hookrightarrow X$? Specialisations
across a one-parameter family of cycle embeddings?

\emph{(ii)} Morphisms of reference curves: holomorphic maps
$C \to C'$? \'Etale maps? Rank-preserving embeddings?

\emph{(iii)} Compatibility between $(\Sigma_{d-1}, C)$ data: does
deforming the cycle $\Sigma_{d-1}$ while keeping $C$ fixed produce an
isomorphism of chiral algebras? A natural transformation?

\smallskip
\noindent\textbf{HEAL.} Define the specialisation
$(\infty, 1)$-category $\mathrm{Slab}_d$:

\begin{definition}[The specialisation $(\infty, 1)$-category $\mathrm{Slab}_d$]
\label{def:slab-d}
Fix a smooth complex CY$_d$ variety $X$ with CY form $\Omega_X$. The
category $\mathrm{Slab}_d(X)$ has:

\emph{Objects}: pairs $(\Sigma_{d-1}, C)$ with $\Sigma_{d-1} \subset X$
a real $(d-1)$-dimensional closed cycle (in the sense of smooth
singular homology with the stratified-cycle refinement of
Francis--Gaitsgory 2012, §3.2), and $C \subset X$ a smooth complex
algebraic curve, together with a transversality datum
$\Sigma_{d-1} \pitchfork C$ (equivalently, a choice of tubular
neighbourhood $\Sigma_{d-1} \times [0, \varepsilon) \hookrightarrow X$
intersecting $C$ in a $0$-dimensional subvariety).

\emph{Morphisms}: a $1$-morphism
$(\Sigma, C) \to (\Sigma', C')$ is a pair $(\varphi, \psi)$ with
$\varphi\colon \Sigma \to \Sigma'$ a smooth bordism in $X$ (a
$d$-dimensional submanifold $W \subset X \times [0, 1]$ with
$\partial W = \Sigma \sqcup \Sigma'$) and $\psi\colon C \to C'$ an
holomorphic map of smooth curves over $X$.

\emph{$2$-morphisms}: homotopies of bordisms (holding $C = C'$ fixed)
and homotopies of holomorphic maps.

The disjoint-union symmetric-monoidal structure is
$(\Sigma, C) \otimes (\Sigma', C') = (\Sigma \sqcup \Sigma', C \sqcup C')$,
with unit $(\emptyset, \emptyset)$ and braiding induced by the
$\bS^\bullet$-framing of $X$. $\mathrm{Slab}_d(X)$ is a symmetric
monoidal $(\infty, 1)$-category.
\end{definition}

\emph{Rationale.} ``Slab'' is Dimofte's terminology for a
$2d$-dimensional slab $C \times \Sigma_{d-1}$ obtained by
specialisation of a $2d$-dimensional holomorphic theory along the
$(d-1)$-cycle $\Sigma_{d-1}$ with reference curve $C$.  The chapter
uses slab geometry implicitly in the K3$\times E$ and $\C^3$
computations; making $\mathrm{Slab}_d$ explicit is a conservative
augmentation that resolves ambiguity.

\begin{proposition}[Stage 2 is a symmetric-monoidal 2-functor on
$\mathrm{Slab}_d(X)$]
\label{prop:spch-slab-d-functor}
The Stage 2 specialisation
\[
 \SpCh\colon \EdHolFA(X)^{(\infty,1)} \times \mathrm{Slab}_d(X)
 \longrightarrow E_{n(d)}\text{-}\HolFA(\bar M_d)^{(\infty,1)}
\]
extends Proposition~\ref{prop:spch-infty1-kernel} to a symmetric-monoidal
$(\infty, 2)$-functor in the second argument:
\begin{enumerate}[label=\textup{(\arabic*)}]
\item On objects: $\SpCh_{\Sigma_{d-1}, C}(\cF) = ((\pi_C)_*(\pi_X^* \cF \otimes^{\bL} \cO_{\Sigma_{d-1} \times C}))^{E_{n(d)}}$.
\item On $1$-morphisms: a bordism $W\colon \Sigma \to \Sigma'$ induces a
chiral-algebra morphism $\SpCh_W(\cF)\colon \SpCh_{\Sigma, C}(\cF) \to
\SpCh_{\Sigma', C}(\cF)$ by integration along $W$.
\item On $2$-morphisms: a homotopy of bordisms induces a chain-homotopy
of chiral-algebra morphisms.
\item Symmetric-monoidal: disjoint union of specialisation data
produces tensor-product chiral algebras (K\"unneth on the curve).
\end{enumerate}
\end{proposition}

\emph{Sketch.} Factorisation homology is a symmetric-monoidal functor
$\int_{(-)}\colon \mathrm{Mfld}_d^{\mathrm{fr}} \to
\mathrm{Alg}_{E_0}$ (Lurie HA 5.5.3, Ayala--Francis 2015). Restricting
to the sub-$(\infty, 1)$-category $\mathrm{Slab}_d(X) \subset
\mathrm{Mfld}_d$ of pairs in $X$ and composing with holomorphic
restriction to $C$ yields the desired $(\infty, 2)$-functor.
K\"unneth multiplicativity is the disjoint-union axiom on the source.

\emph{What was silently assumed.} That
$\SpCh_{\Sigma_{d-1}, C}$ depends on $(\Sigma_{d-1}, C)$ ``only up to
homotopy of cycles''.  The correct statement is that it depends on
the \emph{bordism class} of $\Sigma_{d-1}$ in $X$ and on the
\emph{holomorphic isomorphism class} of $C$; deforming within a
bordism class produces \emph{equivalent} chiral algebras (not equal).
The sibling $\mathfrak g_{\Delta_5}$ / $\mathfrak g_{\Delta_5}^{\mathrm{Enr}}$
/ Igusa-$\chi_{10}$ outputs
(Theorem~\ref{thm:g-delta5-is-sp-k3}; Remark
\ref{rem:multiple-e1-specialisations}) are \emph{not} in the same
bordism class (they correspond to different sections of different
fibrations of $K3 \times E$), hence genuinely different objects of
$E_1\text{-}\HolFA(E)$.

% =====================================================================
\section*{Cycle 5 -- ATTACK (e): is $(U3)$ a theorem or a defining property?}
% =====================================================================

\noindent\textbf{Strike.} Property $(U3)$ reads: the native operadic
level of $\Phi_d(\cC)$ is $n(d) = 1 - d$ (so $\infty, 2, 1$ at
$d = 1, 2, {\geq}3$), and for $d \geq 3$ the braided $E_2$ structure
is recovered on the Drinfeld centre $\cZ(\Rep^{E_1}(\Phi_d(\cC)))$.
Which is it: a theorem about $\Phi_d$, or part of the definition?

\smallskip
\emph{If theorem}: then there is an independent definition of $\Phi_d$
from which $(U3)$ follows. In the chapter's construction
(four-step passage + Stage 1 + Stage 2), the native level is
\emph{computed}: Gerstenhaber-bracket-degree $1 - d$ propagates via
Stage 1 ($E_d$-hFA) through Stage 2 (Dunn restriction to $E_{n(d)}$).
\emph{But} the Drinfeld-centre compatibility clause (``braided $E_2$
recovered on $\cZ(\Rep^{E_1}(A_\cC))$ for $d \geq 3$'') is not
automatic: it asserts a specific equivalence of braided monoidal
categories, which is a theorem about the
Ben-Zvi--Francis--Nadler centre calculation.

\emph{If defining property}: $(U3)$ is part of the characterising
axioms, and its consistency has to be verified on examples (as
property $(U4)$ requires). This is the weaker reading.

\smallskip
\noindent\textbf{HEAL.}  Split $(U3)$ into two clauses:

\begin{theorem}[$(U3a)$ native operadic level -- theorem]
\label{thm:u3a-native-level}
For every smooth proper CY$_d$ category $\cC$, the Stage~2 output
$\SpCh_{\Sigma_{d-1}, C}(\PhiFA_d(\cC))$ carries a native
$E_{n(d)}$-structure with $n(d) = 1, 2, \infty$ at $d \geq 3, 2, 1$
respectively. This is forced by the Gerstenhaber-bracket-degree
$(1 - d)$ on $\HH^\bullet(\cC)$ and the Dunn-restriction dimension
count $n(d) = d - (d - 1) = 1$ (for $d \geq 3$) coming from
factorisation homology over the $(d-1)$-cycle $\Sigma_{d-1}$.
\end{theorem}

\begin{theorem}[$(U3b)$ Drinfeld-centre compatibility -- theorem at $d \leq 3$, conjecture at $d \geq 4$]
\label{thm:u3b-drinfeld-centre}
For $d \geq 3$, the Ben-Zvi--Francis--Nadler Drinfeld-centre
identification
\[
 \cZ\bigl(\Rep^{E_1}(\Phi_d(\cC))\bigr)
 \;\simeq\;
 \Rep^{E_2}\bigl(\PhiFA_d(\cC)\bigr)^{\mathrm{centred}}
\]
(in the sense that the $E_2$-braided monoidal category on the right
is the centralised half-braiding category, Ben-Zvi--Francis--Nadler
\emph{Integral transforms and Drinfeld centers in derived algebraic
geometry}, J.\ Amer.\ Math.\ Soc.\ 23 (2010) Prop.\ 5.5.1).
\emph{Status}: proved at $d = 3$ on toric CY$_3$ inputs (via the
Arakawa--Kapranov--Vasserot quantum groups), conjectural at $d \geq 4$
(the iterated Drinfeld-centre picture requires $E_n$ Morita
equivalences at $n \geq 3$, open in general).
\end{theorem}

\emph{Correction.} $(U3)$ as stated in the G4 specification and in the
platonic-ideal memory conflates two statements of different epistemic
status. $(U3a)$ is a theorem; $(U3b)$ is a theorem at $d = 3$ on a
restricted locus and conjectural in general. The manuscript should
split them.

% =====================================================================
\section*{Cycle 6 -- ATTACK (f): is the native-level stratification
$n(d) \in \{\infty, 2, 1\}$ forced?}
% =====================================================================

\noindent\textbf{Strike.} The stratification $n(d) = \infty, 2, 1$ at
$d = 1, 2, {\geq}3$ is stated as the defining feature of the
target category (Proposition~\ref{prop:native-en-level}).  Two
distinct questions:

\smallskip
\emph{(i)} \textbf{Why does $n(d)$ not vary continuously with $d$?}
One might expect $n(d) = d$ or $n(d) = 2d - 1$ or $n(d) = \max(1, 4-d)$,
any of which would produce different stratification. Why the exact
$\{\infty, 2, 1\}$?

\emph{(ii)} \textbf{Is the $n(d = 1) = \infty$ really forced?}
The Gerstenhaber bracket on $\HH^\bullet(\Coh(E))$ is of degree $0$
(since $d = 1$), and a bracket of degree $0$ on a graded-commutative
algebra is a Poisson bracket, which upgrades
$\HH^\bullet(\Coh(E))$ to an $E_2$-algebra, not an $E_\infty$-algebra.
The stated $n(d = 1) = \infty$ is the \emph{image} in chiral algebras
on a curve (where the $E_2$ structure becomes $E_\infty$ because the
bracket vanishes on the 1-dimensional abelian Lie conformal algebra
$\frakL_{\Coh(E)}$).

\smallskip
\noindent\textbf{HEAL.}  Precise derivation of $n(d)$:

\begin{proposition}[Native operadic level from Gerstenhaber degree]
\label{prop:native-en-heal}
Stage 1 output $\PhiFA_d(\cC) \in \EdHolFA(X)$ is
natively $E_d$-algebraic on $X$ (Kontsevich formality).  Factorisation
homology over a closed $(d-1)$-dimensional manifold
$\Sigma_{d-1}$ induces Dunn restriction
\[
 \int_{\Sigma_{d-1}}\colon \EdHolFA(X) \to E_{d - (d-1)}\text{-}\HolFA(X \setminus \Sigma_{d-1}) = \Eone\text{-}\HolFA(X \setminus \Sigma_{d-1})
\]
for $d \geq 2$, and restriction to a curve $C$ produces
$E_{n(d)}\text{-}\HolFA(C)$ with
\[
 n(d) \;=\; \begin{cases}
 \infty & d = 1\; (\text{Gerstenhaber bracket vanishes,\ Lie conformal algebra abelian}) \\
 2 & d = 2\; (\text{Kontsevich--Vlassopoulos\ $\bS^2$-framing reintroduces the\ $E_2$}) \\
 1 & d \geq 3\; (\text{Dunn restriction\ $E_d \to E_1$,\ no reintroduction})
 \end{cases}
\]
The discontinuity at $d = 2$ is the Kontsevich--Vlassopoulos
$\bS^2$-framing: it is a \emph{chain-level} addition to the native
Dunn restriction, unique to $d = 2$ (at $d \geq 3$ the analogous
$\bS^d$-framing exists but does not upgrade $E_1$ back to $E_2$ on the
curve). The $\infty$ at $d = 1$ is because the Schouten--Nijenhuis
bracket vanishes on $\HH^\bullet(\Coh(E))$ by
$\PV^1(E) = H^0(T_E) = 0$ (elliptic curve has no global vector
fields modulo translations, and modulo translations the bracket
vanishes).
\end{proposition}

\emph{Correction.} The stratification is \emph{forced by geometry}, not
conventional: the Schouten--Nijenhuis bracket on
$\HH^\bullet(D^b(\Coh(E)))$ vanishes because $\PV^1(E)/\text{translations} = 0$;
the $\bS^2$-framing at $d = 2$ is a $d$-specific Kontsevich--Vlassopoulos
datum; at $d \geq 3$ Dunn restriction to $E_1$ is pure Lurie-HA 5.1.2.2
dimension arithmetic.

\smallskip
\emph{Silent assumption now explicit.}  The claim
$n(d = 1) = \infty$ depends on the specific CY$_1$ input
$\Coh(E)$ (or its semi-simple variants). For non-CY curve categories
(e.g., $D^b(\Coh(\bP^1))$, which is a non-CY Fano category), the
output is not an $E_\infty$-chiral algebra. This is why the source of
$\Phi_1$ is restricted to $\CY_1\text{-}\mathrm{Cat}^{\mathrm{dg}}$:
elliptic curves and their finite quotients.

% =====================================================================
\section*{Cycle 7 -- ATTACK (g): do ``many BKMs from one CY$_3$'' really
come from two-stage factorisation, or is this sleight of hand?}
% =====================================================================

\noindent\textbf{Strike.} Corollary~\ref{wn:cor:plat-many-bkms}
(\texttt{platonic\_synthesis\_waves\_11\_through\_16.tex}): ``Borcherds
Monster and Igusa $\mathfrak g_{\Delta_5}$ are sibling $E_1$-specialisations
of a common $E_3$-hFA framework.'' The two-stage factorisation
$\Phi_3 = \SpCh_{\Sigma_2, C} \circ \PhiFA_3$ allegedly explains this
as: one Stage-1 output $\PhiFA_3(\cC)$, multiple Stage-2
specialisations $\SpCh_{\Sigma_2, C}$. Is this a theorem or a
restructured language-game?

\smallskip
\emph{Concern}: The Monster $V^\natural$ is constructed from the Leech
lattice and $\mathrm{II}_{25,1}$, of rank 26. Igusa's $\mathfrak{g}_{\Delta_5}$
comes from $K3 \times E$ with $\Sigma_2 = K3$ and $C = E$, of rank 3.
These are \emph{different ranks, different lattices, different BKMs}.
In what sense are they ``siblings''?

\emph{Sub-concern (a)}: Is there a single CY$_3$ category $\cC$ whose
$\PhiFA_3(\cC)$ specialises under distinct $\Sigma_2$ to both $V^\natural$
and $\mathfrak{g}_{\Delta_5}$?  No -- the Monster requires the Leech
lattice (rank 24) whereas $\mathfrak g_{\Delta_5}$ uses
$\Lambda^{2,1}_{II}$ (rank 3). They correspond to different CY$_3$
inputs.

\emph{Sub-concern (b)}: What is true: there is a family of compact
CY$_3$ inputs, each admitting multiple $\Sigma_2$-fibrations, each
producing multiple $\Eone$-chiral-algebra outputs. For $K3 \times E$
the five CHL siblings $\Phi_N$, $N \in \{1, 2, 3, 4, 6\}$, come from
five distinct $\Sigma_2^{(N)}$-cycles on the same input
(Theorem~\ref{thm:g-delta5-is-sp-k3}, Remark
~\ref{rem:multiple-e1-specialisations},
\texttt{chapters/theory/cy\_to\_chiral.tex}:482).

\smallskip
\noindent\textbf{HEAL.} Split the claim into two:

\begin{theorem}[``Multiple siblings from one CY$_3$'' -- $K3 \times E$]
\label{thm:multi-siblings-k3e}
Fix $X = K3 \times E$ and $\cC = D^b(\Coh(X)) \in \CY_3\text{-}\mathrm{Cat}^{\mathrm{dg}}$.
For each $N \in \{1, 2, 3, 4, 6\}$ there is a specialisation cycle
$\Sigma_2^{(N)} \subset X$ and a reference curve $C = E$ such that
$\SpCh_{\Sigma_2^{(N)}, E}(\PhiFA_3(\cC))$ is the vertex envelope
$U_{\mathrm{ch}}(\mathfrak g^{(N)})$ of a BKM superalgebra with
Macdonald denominator $\Phi_N$ of weight $\kBKM(\Phi_N) = c_N(0)/2$.
The five outputs are distinct chiral algebras on $E$, sibling
$\Eone$-shadows of the single canonical
$\PhiFA_3(\cC) \in E_3\text{-}\HolFA(X)$.
\end{theorem}

\emph{Proof sketch.} The CHL orbifolds
$(K3 \times E)/\Z_N$ are classified by Nikulin's symplectic
automorphisms of K3 of orders 1, 2, 3, 4, 6. Each $\Z_N$-orbifold
gives a distinct $\Sigma_2^{(N)}$ (the section-times-invariant-circle
of the $\Z_N$-action). The $\SpCh_{\Sigma_2^{(N)}, E}$-output is the
Oberdieck--Pixton CHL partition function
$Z^{CHL_N}_{\mathrm{DT}} = C \cdot \Phi_N^{-2}$, which by
Borcherds--Weyl rigidity (Borcherds 1998) is the character of
$U_{\mathrm{ch}}(\mathfrak g^{(N)})$ with weight
$\kBKM = c_N(0)/2$. The claim that the five outputs are siblings of a
common Stage-1 canonical $\PhiFA_3(\cC)$ is the content of the
two-stage factorisation: Stage-1 produces the canonical $E_3$-hFA on
$X$, and Stage-2 provides the five different specialisations.

\begin{corollary}[Monster $\mathfrak g_{V^\natural}$ is NOT a sibling of $\mathfrak g_{\Delta_5}$ via $\Phi_3$]
\label{cor:monster-not-sibling}
The Monster Lie algebra $\mathfrak{g}_{V^\natural}$ (constructed on
$\mathrm{II}_{25,1}$) does \emph{not} arise as a specialisation of
$\PhiFA_3(D^b(\Coh(K3 \times E)))$ via any $\Sigma_2$. Witness: the
Monster root lattice has rank 26 and uses the Leech lattice of rank
24; $K3$ has second Betti number 22 = rank$(U^3 \oplus E_8(-1)^2)$;
no CY$_3$ Hochschild lattice of rank 22 surjects onto the Leech
lattice of rank 24. The Monster is obtained from a distinct CY$_3$
input -- conjecturally a Borcherds--Nikulin $\mathrm{II}_{25,1}$-type
lift -- which does not lie in the $K3 \times E$-sibling family.
\end{corollary}

\emph{Correction.} The language ``many BKMs from one CY$_3$'' is
literal for CHL $\Phi_N$-siblings on $K3 \times E$ (five outputs, one
input, five cycles); it is \emph{not} literal for Monster vs.\
Igusa (distinct CY$_3$ inputs, not shared Stage-1). The
\texttt{platonic\_synthesis\_waves} corollary should be narrowed to
``five CHL siblings from one $K3 \times E$ input'' to be a theorem.

% =====================================================================
\section*{HEAL -- worked rigorous constructions}
% =====================================================================

\subsection*{H1. $\Phi_2$ on $D^b(\Coh(K3))$: the Mukai-Heisenberg,
chain-level}

\noindent The rigorous chain-level construction of
$\Phi_2(D^b(\Coh(K3))) = \cH_{\mathrm{Muk}}$ (rank-24 Heisenberg on the
Mukai lattice) proceeds in four steps. Full detail at
\texttt{cy\_to\_chiral.tex}:2012--2143; here the bones.

\smallskip
\textbf{Step 1 (HKR -- Hochschild homology).} For $S$ a smooth
projective K3, HKR gives
$\HH_n(D^b(\Coh(S))) = \bigoplus_{q - p = n} H^q(S, \Omega^{2-p}_S)$.
K3 Hodge diamond: $h^{0,0} = h^{0,2} = 1$, $h^{1,1} = 20$, all else
zero. Hence $\HH_{-2} = \C$, $\HH_0 = \C^{22}$, $\HH_2 = \C$; odd
groups vanish. Total: $\dim \HH_\bullet = 24 = \dim H^*(S, \C)$.

\textbf{Step 2 (Gerstenhaber bracket vanishes).} Under HKR,
$\HH^\bullet(\cC) = \PV^\bullet(S)$ with Schouten--Nijenhuis bracket.
$\PV^0(S) = H^0(\cO_S) = \C$, $\PV^1(S) = H^0(T_S) = 0$ (K3 has no
vector fields), $\PV^2(S) = H^0(\wedge^2 T_S) = H^0(\cO_S) = \C$ via
the CY-form isomorphism $\wedge^2 T_S \cong \cO_S$. All
Schouten--Nijenhuis brackets vanish: $\{\PV^0, \PV^0\} \subset \PV^{-1} = 0$,
$\{\PV^0, \PV^2\} \subset \PV^1 = 0$, $\{\PV^2, \PV^2\} \subset \PV^3 = 0$.
Hence the Lie conformal algebra $\frakL_\cC = \HH^\bullet(\cC)$ is
abelian: $[a_\lambda b] = 0$.

\textbf{Step 3 ($\bS^2$-framing upgrades to $E_2$).} Kontsevich--Vlassopoulos
2022 provides the pre-CY structure: every cyclic $\Ainf$ structure of
dimension 2 gives an $\bS^2$-framing on $\HH_\bullet$. For K3 this is
explicit: $\HH_\bullet(D^b(\Coh(S))) = H^*(S, \C)$ as a graded
vector space with the $\bS^2$-framing coming from Poincar\'e
duality on the surface.

\textbf{Step 4 (Mukai pairing from Hochschild trace).} The CY trace
$\mathrm{tr}\colon \HH_2(\cC) \to \C$ is integration on the top class;
the Markarian--C\u{a}ld\u{a}raru theorem identifies the induced
pairing on $\HH_\bullet = H^*(S, \C)$ as the Mukai pairing
$\omega_{\mathrm{Muk}}(\alpha, \beta) = \int_S \alpha \cup \beta \cup \mathrm{td}^{1/2}(S)$
with $\mathrm{td}^{1/2}(S) = 1 + [\mathrm{pt}]$. Explicitly:
\[
 \omega_{\mathrm{Muk}} = \begin{pmatrix} 0 & 0 & -1 \\ 0 & Q_{22} & 0 \\ -1 & 0 & 0 \end{pmatrix},
 \quad Q_{22} = U^3 \oplus E_8(-1)^2, \quad \text{signature}\ (4, 20).
\]

\textbf{Output.} Since $\frakL_\cC$ is abelian, the factorisation
envelope is the free Heisenberg: $\Phi_2(D^b(\Coh(S))) = \Fact_C(\frakL_\cC) = \mathrm{Heis}(H^*(S, \C), \omega_{\mathrm{Muk}})$
on the reference curve $C$. Character $= 1/\eta^{24}$; modular
characteristic $\kch = 2 = \chi(\cO_S)$; shadow class $\mathrm G$
(Gaussian, $m_k = 0$ for $k \geq 3$); $E_2$-enhancement from
$\bS^2$-framing upgrades automatically to $E_\infty$ on the abelian
Heisenberg (the free Heisenberg is $E_\infty$ because commutative
factorisation algebras are $E_\infty$).

\smallskip\noindent\emph{Status:} proved chain-level (Theorem
\ref{thm:phi-k3-explicit}, \texttt{cy\_to\_chiral.tex}:1994).
Verified by \texttt{phi\_k3\_explicit\_evaluation.py} (55 tests).
Functoriality on morphisms (property $(U2)$): tested on the Mukai
transform $\mathrm{K3} \to \mathrm{K3}$, chain-level compatibility
pending. Independent verifications: Kummer route (Conjecture
\ref{conj:kummer-route}); $K3^{[2]}$ Hilbert scheme (Nakajima
factorisation homology); heterotic $\sigma$-model at K3 point.

\subsection*{H2. $\Phi_3$ on toric $\C^3$: $\CoHA(\C^3) \mapsto Y^+(\widehat\fgl_1)$}

\noindent Following Schiffmann--Vasserot 2013 (arXiv:1202.2756) and the
two-stage factorisation at the point-specialisation
$(\Sigma_2, C) = (\mathrm{pt}, \C)$:

\textbf{Step 1 (input).} $\cC = \CoHA(\C^3)$, the critical CoHA of the
quiver-with-potential $(Q^{(3)}, W^{(3)})$ with one vertex and three
loops $x, y, z$ and potential $W = x[y, z]$ (Kontsevich--Soibelman
2011 arXiv:1006.2706 §8). As a $\CY_3$-dg-category, $\cC$ is smooth
proper with $\HH^0(\cC) = \C$ (rank-one centre) and non-degenerate
cyclic pairing of degree $-3$ from the Kontsevich--Soibelman DT-trace.

\textbf{Step 2 ($\PhiFA_3$ via 6d hCS).} Proposition
\ref{prop:phifa-infty1-kernel} realises $\PhiFA_3(\cC)$ as the
quantum BV factorisation algebra
$\Obs^{\mathrm{hCS}}(\C^3, \widehat\fgl_1)$ on $\C^3$ with gauge
dg-Lie $\widehat\fgl_1 = \HH^\bullet(\cC) = \C[x, y, z]/(\text{Jacobi})$
(the Jacobi algebra of the potential). By Costello--Li 2020, this
$E_3$-holomorphic factorisation algebra is the Costello--Yagi
realisation of the affine Yangian (arXiv:1306.4657).

\textbf{Step 3 ($\SpCh_{\mathrm{pt}, \C}$ at the point-specialisation).}
Factorisation homology over the point
$\int_{\mathrm{pt}} = \mathrm{id}$ (trivial integration) followed by
restriction to the reference curve $C = \C$ gives
\[
 \SpCh_{\mathrm{pt}, \C}(\PhiFA_3(\cC)) \;=\; \Obs^{\mathrm{hCS}}(\C^3, \widehat\fgl_1)\big|_\C \;=\; Y^+(\widehat\fgl_1),
\]
the positive half of the affine Yangian of $\widehat\fgl_1$.

\smallskip\noindent\emph{Output.} $\Phi_3(\CoHA(\C^3)) = Y^+(\widehat\fgl_1)$
as $\Eone$-chiral algebra on $\C$. The Drinfeld-centre doubling
(property $(U3b)$) applied to $Y^+$ recovers the full Yangian
$Y(\widehat\fgl_1)$; the evaluation chain
$Y^+(\widehat\fgl_1) \hookrightarrow Y(\widehat\fgl_1) \xrightarrow{\mathrm{ev}_\lambda} \mathrm{End}(\mathcal W_{1+\infty}[\lambda])$
exhibits $\mathcal W_{1+\infty}$ at $c = 1$ as a module (not the
$\Phi_3$-output itself).

\smallskip\noindent\emph{Status.} Stage 1 proved (Costello--Yagi 2016
arXiv:1306.4657; Costello--Gaiotto 2018 arXiv:1804.09231). Stage 2
proved at point-specialisation (Schiffmann--Vasserot 2013). Chain-level
identification $\PhiFA_3(\cC)\big|_\C = Y^+$ proved under vanishing of
the rank-3 Casimir (simply-laced $\widehat\fgl_1$ has $d^{abc} = 0$).
Independent verifications: Arakawa--Kuwabara--Thielmans 2023 (twisted
module category of $\mathcal W_{1+\infty}$); Oblomkov--Rozansky 2019
(Hilbert scheme cohomology of $\C^3$); Gaiotto--Rap\v{c}\'ak 2020 (6d hCS
partition function on $\C^3$).

\subsection*{H3. $\Phi_3$ on $K3 \times E$: the Oberdieck--Pixton
identity}

\noindent On $X = K3 \times E$ with canonical cycle $\Sigma_2 = K3$
(the K3-fibre of $p_{K3}: X \to K3$) and reference curve $C = E$:

\smallskip
\textbf{Step 1 (input).} $\cC = D^b(\Coh(K3 \times E))$. By Bondal--Orlov,
$\cC = \Perf^{\mathrm{dg}}(K3 \times E) \simeq \Perf^{\mathrm{dg}}(K3) \boxtimes \Perf^{\mathrm{dg}}(E)$.
Each factor is CY (K3 is strict CY$_2$, $E$ is CY$_1$ via
$\omega_E = \cO_E$). The total category is $\CY_3$-dg with
$\HH^0 = \C$, cyclic pairing in degree $-3$ from the product trace.

\textbf{Step 2 ($\PhiFA_3$ via K\"unneth).} By
Proposition~\ref{prop:phi-d-kunneth-split-product}, Stage~1 factorises:
\[
 \PhiFA_3(D^b(\Coh(K3 \times E))) \;=\; \PhiFA_2(D^b(\Coh(K3))) \boxtimes \PhiFA_1(\Coh(E)),
\]
a product of two factorisation algebras on $K3 \times E$. The K3 factor
is the canonical $E_2$-hFA underlying the rank-24 Mukai-Heisenberg
(see H1 above); the $E$ factor is the canonical $E_1$-hFA underlying
the rank-2 Heisenberg $H_1 \otimes H_2 = \mathrm{Heis}(H^0(E) \oplus H^1(E))$.

\textbf{Step 3 ($\SpCh_{K3, E}$).} Factorisation homology over
$\Sigma_2 = K3$ (the K3-fibre) integrates out the $K3$ factor and
leaves the $E$ factor on the reference curve $E$:
\[
 \SpCh_{K3, E}(\PhiFA_3(\cC)) \;=\; \int_{K3} \PhiFA_2(D^b(\Coh(K3))) \;\boxtimes\; \PhiFA_1(\Coh(E))\big|_E.
\]

The integral $\int_{K3}$ is the K3 factorisation-homology trace,
identified by Oberdieck--Pixton 2017 (arXiv:1706.10100 Thm 3) with the
Igusa cusp form:
\[
 Z^{K3 \times E}_{\mathrm{DT}}(q, y, Q) \;=\; C \cdot \Delta_5(Z)^{-2} \big|_{(q, y, Q)}.
\]
The character of the $\Eone$-chiral algebra output is
$\Delta_5(Z)^{-1}\big|_E$ (the BPS partition function on the section of
the elliptic fibration), and by Borcherds--Weyl rigidity 1998, this
character is uniquely realised by the vertex envelope
$U_{\mathrm{ch}}(\mathfrak{g}_{\Delta_5})$ of the Igusa BKM superalgebra.

\textbf{Output.} $\Phi_3^{(K3, E)}(D^b(\Coh(K3 \times E))) = U_{\mathrm{ch}}(\mathfrak{g}_{\Delta_5})$
on the curve $E$. Weight $\kBKM = c_{\phi_{0,1}}(0, 0)/2 = 10/2 = 5$.

\smallskip\noindent\emph{Status.} Stage 1 chain-level via K\"unneth
(Proposition~\ref{prop:phi-d-kunneth-split-product}, unconditional on
split products of CY factors).  Stage 2 chain-level at
$(\Sigma_2 = K3, C = E)$ via Oberdieck--Pixton
(arXiv:1706.10100 Thm 3). Theorem
\ref{thm:g-delta5-is-sp-k3} inscribed. Independent verifications:
Gritsenko 1999 $\Delta_5$ Fourier coefficients; DMZ denominator
identity; Mathieu moonshine for the mock-Jacobi form attached to the
K3 anti-symplectic involution.

\emph{Other $(\Sigma_2, C)$ choices on $K3 \times E$}: the CHL
siblings $\Phi_N$, $N \in \{2, 3, 4, 6\}$, come from symplectic
K3-automorphisms of order $N$ (Nikulin's classification) and give
distinct outputs
$U_{\mathrm{ch}}(\mathfrak{g}^{(N)})$ with weights
$\kBKM(\Phi_N) = c_N(0)/2 \in \{4, 3, 2, 1\}$.
Theorem~\ref{thm:multi-siblings-k3e} above.

% =====================================================================
\section*{Opens and harmonies}
% =====================================================================

\subsection*{Opens (explicit, after this audit)}

\begin{enumerate}[leftmargin=*]
\item \textbf{Uniqueness off the Koszul-self-dual locus.} Theorem
\ref{thm:phi-uniqueness-sharp} restricts uniqueness to
$\cC_d^{\mathrm{Kosz}}$.  Off this locus, a Fourier--Mukai
rigidification of $\Phi_d$ via $\cK_{\PhiFA_d}$ is needed. Open.

\item \textbf{Functoriality on morphisms ($(U2)$) at $d \geq 2$.}
Conjecture~\ref{conj:phi-d-functoriality} is open at $d \geq 2$
(chain-level); the Mukai transform is the first test case at $d = 2$.

\item \textbf{Chain-level $E_d$-formality of $\HH^\bullet(\cC)$ at
$d \geq 4$.} Operadic formality of $E_d$ is Fresse Vol~I Thm
14.1.A; category-level formality of a general smooth proper CY$_d$
Hochschild complex is open at $d \geq 4$.

\item \textbf{BV-locality $\mathrm{Obs}_{\mathrm{BV}}$ at $d \geq 3$.}
Perturbative proof in Costello--Li 2020; non-perturbative convergence
is the last remaining obstruction at $d = 3$
(Proposition~\ref{prop:hopf-fibration-decomposition}).

\item \textbf{$(U3b)$ Drinfeld-centre compatibility at $d \geq 4$.}
Theorem~\ref{thm:u3b-drinfeld-centre} is a theorem at $d = 3$ on toric;
at $d \geq 4$ requires iterated $E_n$ Morita equivalences for
$n \geq 3$, open in general.

\item \textbf{The general ``multi-siblings'' for non-$K3 \times E$
CY$_3$.} Theorem~\ref{thm:multi-siblings-k3e} covers CHL siblings on
$K3 \times E$. Analogous theorems for Enriques $\times E$ orbifolds,
for Humbert surfaces, for Niemeier-class inputs, are open (conjectured
in \texttt{cy\_to\_chiral.tex}:482, last sentence).

\item \textbf{$\mathrm{Slab}_d$ as a self-contained $(\infty, 2)$-category.}
Definition~\ref{def:slab-d} above needs verification of the
$(\infty, 2)$-categorical axioms (associator, pentagon, unit constraints).
Open (sketch only).
\end{enumerate}

\subsection*{Harmonies with other wave-17 G-identities}

\begin{itemize}[leftmargin=*]
\item \textbf{G1 (Koszul reflection)}: Vol~I Theorem A gives the
bar-cobar adjunction on the Koszul locus. Theorem
\ref{thm:phi-uniqueness-sharp} above restricts the uniqueness of
$\Phi_d$ to this locus; $\Phi_d$ on its domain commutes with
bar-cobar on its target.

\item \textbf{G2 (Arnold connection)}: Vol~I climax -- the Gaitsgory--Kazhdan
KZ/KZB connection on $\mathrm{Conf}_n$. The K3$\times E$ shadow of
$\Phi_3$ restricts to a vertex operator realisation of an elliptic
Arnold connection on $\mathrm{Conf}_n(E)$.

\item \textbf{G3 (Universal holography)}: Vol~II climax -- boundary
$E_2$-chiral, bulk $E_3$-top, derived centre. The Drinfeld-centre clause
$(U3b)$ of $\Phi_d$ is the Vol~III-specific instance of the
$E_2 = \cZ(E_1)$ holographic identity: CY$_d$ input produces
$\Eone$-chiral boundary, whose Drinfeld centre is the $E_2$ bulk.

\item \textbf{Universal trace identity (UTI)}: $K \circ \Phi = \kappa_{\mathrm{ch}}$
at the level of functors; on K3-fibered CY$_3$ $(K \circ \Phi)/2 = \kBKM$.
The sharpening of $\Phi_3$ here strengthens the UTI: the trace
identity holds strictly on Koszul-self-dual $\cC$ and
up-to-central-charge on general $\cC$.

\item \textbf{Six routes to $G(K3 \times E)$ (Vol III Ch
\texttt{cy-c-six-routes-convergence})}: the two-stage factorisation of
this note identifies the six routes as six distinct $(\Sigma_2, C)$
choices on a common Stage~1 output, not as six applications of
$\Phi_3$. This is the content of
Theorem~\ref{thm:multi-siblings-k3e}.
\end{itemize}

% =====================================================================
\section*{Cycle log}
% =====================================================================

\begin{enumerate}[leftmargin=*]
\item[\textbf{C1.}] Uniqueness via $(U1), (U3), (U4)$ is a theorem
\emph{only on the Koszul-self-dual locus}
(Theorem~\ref{thm:phi-uniqueness-sharp}). Counter-example sketch: a
central-charge-twisted $\widetilde\Phi_3$ agrees with $\Phi_3$ on four
canonical inputs but differs on non-canonical compact CY$_3$ with
$h^{0,2} \neq 0$. Fix: restrict to $\cC_d^{\mathrm{Kosz}}$, or
rigidify off by Fourier--Mukai kernel.
\item[\textbf{C2.}] $\Phi_2$ is unconditional \emph{chain-level on
objects} and conjectural \emph{chain-level on morphisms}. The phrase
``$\Phi_2$ is PROVED'' should specify: proved on objects, tested on
the Mukai transform for morphisms. Three inputs (KV $\bS^2$-framing,
Kontsevich $E_2$-formality, CGL locality) all chain-level
unconditional at $d = 2$.
\item[\textbf{C3.}] $E_d$-formality at $d \geq 3$ is operadically proved
(Fresse Vol~I Thm 14.1.A) but \emph{category-level} formality of
$\HH^\bullet(\cC)$ is (a) chain-level proved at $d = 2$, (b)
$(\infty,1)$-proved up-to-contractible-choice at $d = 3$, (c) open
at $d \geq 4$. Silent assumption in the G4 specification: category-level
formality. Fix: distinguish operad vs category.
\item[\textbf{C4.}] $\mathrm{Slab}_d(X)$ defined
(Definition~\ref{def:slab-d}): objects are $(\Sigma_{d-1}, C)$ pairs
with transversality, $1$-morphisms are bordisms + holomorphic maps,
$2$-morphisms are homotopies. Stage 2 is a symmetric-monoidal
$(\infty, 2)$-functor on $\mathrm{Slab}_d$
(Proposition~\ref{prop:spch-slab-d-functor}).
Silent assumption: $\SpCh$ depends on $(\Sigma_{d-1}, C)$ only up to
homotopy of cycles; correct: bordism class + holomorphic-iso class.
\item[\textbf{C5.}] $(U3)$ splits into $(U3a)$ native level (theorem,
Theorem~\ref{thm:u3a-native-level}) and $(U3b)$ Drinfeld-centre
compatibility (theorem at $d = 3$ on toric, conjecture at
$d \geq 4$, Theorem~\ref{thm:u3b-drinfeld-centre}).
\item[\textbf{C6.}] Native-level stratification $\{\infty, 2, 1\}$ is
\emph{forced by geometry}: Schouten--Nijenhuis vanishing at $d = 1$,
KV $\bS^2$-framing at $d = 2$, Dunn restriction $E_d \to E_1$ at
$d \geq 3$. Not a convention; derived in
Proposition~\ref{prop:native-en-heal}.
\item[\textbf{C7.}] ``Many BKMs from one CY$_3$'' is
\emph{literal for CHL siblings on $K3 \times E$} (five
$(\Sigma_2^{(N)}, E)$-choices on one Stage-1 output,
Theorem~\ref{thm:multi-siblings-k3e}), \emph{not literal} for
Monster vs Igusa (distinct CY$_3$ inputs,
Corollary~\ref{cor:monster-not-sibling}).
\end{enumerate}

% =====================================================================
\section*{Healed outputs -- consolidated}
% =====================================================================

\begin{enumerate}[leftmargin=*]
\item[\textbf{H1.}] $\Phi_2(D^b(\Coh(K3))) = \cH_{\mathrm{Muk}}$
(rank-24 Heisenberg at the Mukai pairing) is chain-level proved via
the four-step cyclic-to-chiral passage: HKR $\to$ Schouten--Nijenhuis
vanishing $\to$ $\bS^2$-framing $\to$ Markarian--C\u{a}ld\u{a}raru trace
identification. Signature $(4, 20)$, $\kch = 2$, class G, shadow
depth 0. Unconditional; tested in \texttt{phi\_k3\_explicit\_evaluation.py}.

\item[\textbf{H2.}] $\Phi_3(\CoHA(\C^3)) = Y^+(\widehat\fgl_1)$ at the
point-specialisation $(\Sigma_2, C) = (\mathrm{pt}, \C)$ is chain-level
proved via 6d hCS BV (Costello--Li--Yagi) + Schiffmann--Vasserot 2013.
The full $\mathcal W_{1+\infty}$ at $c = 1$ is obtained only after
Drinfeld-centre doubling (property $(U3b)$). Positive-half vs
full-Yangian distinction is structural.

\item[\textbf{H3.}] $\Phi_3(D^b(\Coh(K3 \times E))) = U_{\mathrm{ch}}(\mathfrak{g}_{\Delta_5})$
at $(\Sigma_2, C) = (K3, E)$ is proved via K\"unneth (Stage 1) +
Oberdieck--Pixton DT/Igusa identity (Stage 2) + Borcherds--Weyl
rigidity. Weight $\kBKM = 5$. The five CHL siblings $\Phi_N$ are
sibling specialisations on the same Stage 1 output; the Monster is
\emph{not} a sibling.

\item[\textbf{H4.}] Uniqueness of $\Phi_d$ holds on
$\cC_d^{\mathrm{Kosz}}$ (Theorem~\ref{thm:phi-uniqueness-sharp}); off
that locus, Fourier--Mukai rigidification is the natural extension.

\item[\textbf{H5.}] $\mathrm{Slab}_d(X)$ is the specialisation
$(\infty, 2)$-category (Definition~\ref{def:slab-d}); Stage 2 is a
symmetric-monoidal $(\infty, 2)$-functor on it
(Proposition~\ref{prop:spch-slab-d-functor}).

\item[\textbf{H6.}] $(U3)$ splits as $(U3a)$ native level (theorem) +
$(U3b)$ Drinfeld-centre (theorem at $d = 3$, conjecture at
$d \geq 4$).  Native-level stratification $\{\infty, 2, 1\}$ is forced
by geometry, not convention.
\end{enumerate}

% =====================================================================
\section*{Hidden assumptions made explicit}
% =====================================================================

\begin{enumerate}[leftmargin=*]
\item \textbf{Koszul-self-duality} was silently used in the uniqueness
argument; now explicit in Theorem~\ref{thm:phi-uniqueness-sharp}.

\item \textbf{Category-level (not just operad-level) $E_d$-formality}
was silently conflated with operad-level; now split by Cycle 3.

\item \textbf{Bordism-class (not homotopy-class) dependence of
$\SpCh_{\Sigma_{d-1}, C}$} was silently assumed; now explicit
in Definition~\ref{def:slab-d}.

\item \textbf{$\PV^1(E) / \text{translations} = 0$} (not just
$\PV^1(E) = \C$) drives the $n(d = 1) = \infty$ stratification; the
translation quotient was silent.

\item \textbf{$\Phi_3$ on non-$K3$-fibred inputs (quintic, local $\bP^2$)}:
the output is \emph{not} a Borcherds lift (the Mukai-lattice
unimodularity fails: $b_3(Q_5) = 204$). Silent assumption: that
$\kBKM$ is defined for every $\Phi_3$-output; correct: $\kBKM$ is
defined only on the K3-fibred sublocus, and $\kappa_{\mathrm{BCOV}}$ takes
its role for the quintic.

\item \textbf{Drinfeld-centre compatibility at $d \geq 4$} is conjectural;
the G4 specification did not scope this (implicit ``for $d \geq 3$'').
Fix: $(U3b)$ proved at $d = 3$ on toric, conjecture at $d \geq 4$
(Theorem~\ref{thm:u3b-drinfeld-centre}).
\end{enumerate}

% =====================================================================
\section*{Convergence: why this matters}
% =====================================================================

The two-stage factorisation
$\Phi_d = \SpCh_{\Sigma_{d-1}, C} \circ \PhiFA_d$ was inscribed in
Wave 11 as the unifying structural statement of Vol III. This audit
locates seven precise failure modes and seven precise fixes. On the
healed statement:

\begin{itemize}[leftmargin=*]
\item The chain-level chain from CY$_d$-category to $\Eone$-chiral
algebra on a curve is unconditional at $d \leq 3$ \emph{on Koszul
inputs} and conditional at $d \geq 4$.
\item The multiplicity of BKMs from one CY$_3$ is genuine for CHL
siblings on $K3 \times E$, illusory for Monster.
\item The universal trace identity $K \circ \Phi = \kch$ has a rigorous
footing on $\cC_d^{\mathrm{Kosz}}$.
\item The three-volume decomposition is \emph{forced} by the native-level
stratification ($\{\infty, 2, 1\}$ is not a convention; it is geometry).
\end{itemize}

This document supports inscription of the sharpened Remark
\ref{rem:phi-uniqueness}, the split of $(U3)$ into $(U3a/U3b)$, and
the definition of $\mathrm{Slab}_d$ into the chapter
\texttt{cy\_to\_chiral.tex}. Without these corrections the G4
specification is optimistic; with them, it is true.

\end{document}
